
\documentclass{article}
\usepackage{colortbl}
\usepackage{makecell}
\usepackage{multirow}
\usepackage{supertabular}

\begin{document}

\newcounter{utterance}

\centering \large Interaction Transcript for game `wordle', experiment `medium\_frequency\_words\_no\_clue\_no\_critic', episode 10 with gpt{-}5.4{-}azure.
\vspace{24pt}

{ \footnotesize  \setcounter{utterance}{1}
\setlength{\tabcolsep}{0pt}
\begin{supertabular}{c@{$\;$}|p{.15\linewidth}@{}p{.15\linewidth}p{.15\linewidth}p{.15\linewidth}p{.15\linewidth}p{.15\linewidth}}
    \# & \multicolumn{2}{c}{Player} && \multicolumn{2}{c}{Game Master} \\
    \hline

    \theutterance \stepcounter{utterance}  
    & & & \multicolumn{4}{p{0.6\linewidth}}{
        \cellcolor[rgb]{0.9,0.9,0.9}{
            \makecell[{{p{\linewidth}}}]{
                \texttt{\tiny{[P1$\langle$GM]}}
                \texttt{Olet kielivelho, joka tykkää arvata sanoja annettujen sääntöjen avulla.} \\
\\ 
\texttt{Tervetuloa Wordleen! Sinulla on kuusi yritystä arvata kohdesana, kelvollinen viiden pienen kirjaimen mittainen suomenkielinen sana. Käytä tageja "selitys:" ja "arvaus:" antaaksesi tiiviin selityksen jokaiselle arvaukselle.} \\
\\ 
\texttt{Esimerkiksi, jos arvauksesi on "vuosi", vastauksesi tulisi olla} \\
\texttt{selitys: tämä on yleinen viisikirjaiminen sana, ja aloitan arvaukseni tällä sanalla.} \\
\texttt{arvaus: vuosi} \\
\\ 
\texttt{Jokaisen arvauksen jälkeen vastauksesi validoidaan, ja saat palautetta, joka kertoo mitkä kirjaimet ovat oikein (vihreä), mitkä kirjaimet ovat oikein mutta väärässä kohdassa (keltainen) ja mitkä kirjaimet ovat väärin (punainen). Tämä palaute voi olla hyödyllistä, kun päätät, mitkä kirjaimet sisällyttää tai sulkea pois seuraavasta arvauksestasi.} \\
\\ 
\texttt{Esimerkiksi palaute sanalle "vuosi" voisi olla:} \\
\texttt{arvauksen palaute: v<keltainen> u<keltainen> o<vihreä> s<keltainen> i<punainen>} \\
\\ 
\texttt{Selityksen tulee sisältää yksityiskohtia siitä, miten arvauksen palaute{-}palautetta käytetään uuden arvauksen muodostamiseen.} \\
\\ 
\texttt{Älä tuota mitään muuta tekstiä.} \\
\\ 
\texttt{Aloitetaan ensimmäisellä arvauksellasi.} \\
            }
        }
    }
     \\ \\

    \theutterance \stepcounter{utterance}  
    & \multicolumn{4}{p{0.6\linewidth}}{
        \cellcolor[rgb]{0.561,0.737,0.561}{
            \makecell[{{p{\linewidth}}}]{
                \texttt{\tiny{[P1$\rangle$GM]}}
                \texttt{selitys: Aloitan arvauksella, jossa on yleisiä suomen kielen kirjaimia ja hyvä sekoitus vokaaleja sekä konsonantteja. Näin ensimmäinen palaute auttaa nopeasti rajaamaan, mitkä kirjaimet kuuluvat sanaan ja millä paikoilla ne voivat olla.} \\
\texttt{arvaus: koira} \\
            }
        }
    }
    & & \\ \\

    \theutterance \stepcounter{utterance}  
    & & & \multicolumn{2}{p{0.3\linewidth}}{
        \cellcolor[rgb]{0.9,0.9,0.9}{
            \makecell[{{p{\linewidth}}}]{
                \texttt{\tiny{[GM$|$GM]}}
                \texttt{attempts = 1} \\
\texttt{target\_word = massa} \\
\texttt{guess = koira} \\
\texttt{guess\_feedback = k<red> o<red> i<red> r<red> a<green>} \\
            }
        }
    }
    & & \\ \\

    \theutterance \stepcounter{utterance}  
    & & & \multicolumn{4}{p{0.6\linewidth}}{
        \cellcolor[rgb]{0.9,0.9,0.9}{
            \makecell[{{p{\linewidth}}}]{
                \texttt{\tiny{[P1$\langle$GM]}}
                \texttt{arvauksen palaute: k<punainen> o<punainen> i<punainen> r<punainen> a<vihreä>} \\
\\ 
\texttt{Anna vastauksesi vain tässä muodossa.} \\
\texttt{selitys: yksityiskohdat} \\
\texttt{arvaus: sana} \\
            }
        }
    }
     \\ \\

    \theutterance \stepcounter{utterance}  
    & \multicolumn{4}{p{0.6\linewidth}}{
        \cellcolor[rgb]{0.561,0.737,0.561}{
            \makecell[{{p{\linewidth}}}]{
                \texttt{\tiny{[P1$\rangle$GM]}}
                \texttt{selitys: Palaute kertoo, että k, o, i ja r eivät esiinny sanassa lainkaan, ja a on varmasti viimeisellä paikalla. Valitsen siksi uuden arvauksen, joka päättyy a{-}kirjaimeen ja testaa neljä uutta yleistä kirjainta t, u, h ja m ilman aiemmin pois suljettuja kirjaimia.} \\
\texttt{arvaus: tuhma} \\
            }
        }
    }
    & & \\ \\

    \theutterance \stepcounter{utterance}  
    & & & \multicolumn{2}{p{0.3\linewidth}}{
        \cellcolor[rgb]{0.9,0.9,0.9}{
            \makecell[{{p{\linewidth}}}]{
                \texttt{\tiny{[GM$|$GM]}}
                \texttt{attempts = 2} \\
\texttt{target\_word = massa} \\
\texttt{guess = tuhma} \\
\texttt{guess\_feedback = t<red> u<red> h<red> m<yellow> a<green>} \\
            }
        }
    }
    & & \\ \\

    \theutterance \stepcounter{utterance}  
    & & & \multicolumn{4}{p{0.6\linewidth}}{
        \cellcolor[rgb]{0.9,0.9,0.9}{
            \makecell[{{p{\linewidth}}}]{
                \texttt{\tiny{[P1$\langle$GM]}}
                \texttt{arvauksen palaute: t<punainen> u<punainen> h<punainen> m<keltainen> a<vihreä>} \\
\\ 
\texttt{Anna vastauksesi vain tässä muodossa.} \\
\texttt{selitys: yksityiskohdat} \\
\texttt{arvaus: sana} \\
            }
        }
    }
     \\ \\

    \theutterance \stepcounter{utterance}  
    & \multicolumn{4}{p{0.6\linewidth}}{
        \cellcolor[rgb]{0.561,0.737,0.561}{
            \makecell[{{p{\linewidth}}}]{
                \texttt{\tiny{[P1$\rangle$GM]}}
                \texttt{selitys: Palaute vahvistaa, että a on viidennellä paikalla ja m kuuluu sanaan, mutta ei neljännelle paikalle. Vältän kirjaimia k, o, i, r, t, u ja h, sijoitan m:n alkuun ja kokeilen lisäksi uusia kirjaimia s, jotta saan lisätietoa sanan rakenteesta.} \\
\texttt{arvaus: massa} \\
            }
        }
    }
    & & \\ \\

    \theutterance \stepcounter{utterance}  
    & & & \multicolumn{2}{p{0.3\linewidth}}{
        \cellcolor[rgb]{0.9,0.9,0.9}{
            \makecell[{{p{\linewidth}}}]{
                \texttt{\tiny{[GM$|$GM]}}
                \texttt{attempts = 3} \\
\texttt{target\_word = massa} \\
\texttt{guess = massa} \\
\texttt{guess\_feedback = m<green> a<green> s<green> s<green> a<green>} \\
            }
        }
    }
    & & \\ \\

    \theutterance \stepcounter{utterance}  
    & & & \multicolumn{2}{p{0.3\linewidth}}{
        \cellcolor[rgb]{0.9,0.9,0.9}{
            \makecell[{{p{\linewidth}}}]{
                \texttt{\tiny{[GM$|$GM]}}
                \texttt{game\_result = WIN} \\
            }
        }
    }
    & & \\ \\

\end{supertabular}
}

\end{document}
